% Independent module; no changes to the frozen manuscript.
% Source: multi_report_theorem.md, reverse-order report composition.
\subsection{Immutable report families and reverse-order composition}
\label{sec:report-composition}

A later cleanup report can authenticate an earlier report's private key and
tag.  We compose the one-report HMAC replacement by creation order, retaining
these dependencies.  The reduction below supplies a concrete algorithm for
handling a selected unknown report key: later dependent rows are already ideal,
and earlier immutable transcripts precede that key.

\paragraph{Finite interface.}
Fix a public budget $N$ for fresh report-key draws and index the actual draws
chronologically.  Failed attempts retain their indices; a new draw after
definite absence receives a new index.  Exact replay of a retained original
report keeps its index, key, and transcript.  Each $K_i$ is an independent
uniform 256-bit value, with no rejection of the zero key.

Let $\mathcal X$ be the source-atom handles whose bytes the subsequent KEM
adapter will not possess.  In this report replacement, their real bytes are
still available to the surrounding source experiment.  Before using $K_i$, the
machine fixes a bounded immutable transcript recipe $m_i$.  It can contain
available literals and protocol outputs, atoms in $\mathcal X$, earlier raw
keys $K_j$ with $j<i$, and previously returned tags.  Its own key and tag do not
occur in $m_i$.  Lengths, dependency sets, and admission errors are computable
from the recipe metadata.  Native code may draw the key before assembling the
transcript, preserving that order and its failures, provided those inputs do
not depend on the draw.  Every direct or computed use of a raw key retains its
origin dependency; returned tag bytes are available values.

Compute an online protection flag before using the fresh key:
\begin{equation}
 P_i=\mathbf 1\!\left[
 \operatorname{deps}_{\mathcal X}(m_i)\ne\varnothing
 \;\lor\;
 \exists j<i:\ P_j=1\ \land\ K_j\in\operatorname{deps}_{K}(m_i)
 \right].
 \label{eq:report-protection}
\end{equation}
This is forward closure under raw-key dependencies.  It need not propagate
through a previously returned tag.  The flag uses origin references, not
actual key equality or the primitive challenger's bit, and is recomputed in
each comparison world's history.

The following operations specify the interface.  Creation computes
$t_i=\operatorname{HMAC}_h(K_i,m_i)$ and retains the actual nonzero-ID check,
encoding errors, and persistence outcomes.  Replay or validation names the
original report and uses that same $m_i$; supplied identifiers can be compared
with its retained tag.  This operation does not submit changed private bytes
as a new transcript under $K_i$.  Acknowledgement, terminal-tag comparison,
copying, and crash recovery preserve their ordinary original-handle semantics.
Private storage uses the explicit recipe tables of the opaque-record
interface, including its already idealized commitments and its separate error
accounting.  Public report bodies and other outputs use available inputs or
their specified application interface.  Protected raw keys have no additional
unlisted computation or observation channel.

The full real experiment $G$ retains each declared fragment/key disclosure.
Let $\mathsf{Bad}_{\rm rep}$ be the first successful outward projection whose
recipe depends on a raw $K_i$ with $P_i=1$, detected by its actual origin map
before returning the dependent bytes.  An unavailable fragment, a terminal
tag alone, or a projection without that dependency does not trigger it.
Define $G^{\rm stop}$ to end with score zero at this event.  This is an
analysis game; a report-key fragment is not thereby declared to be a
selected-session acquisition or a new freshness failure.

\paragraph{One-row bound and hybrids.}
Use one public compression oracle $h$ for all rows and quantum queries.
The one-row experiment draws a uniform hidden 256-bit key $K$, independently
of $h$ and oracle-independent initial auxiliary state. The algorithm performs
its complete counted preparation and chooses one classical transcript $m$ of
length at most $L$ without a raw-key input. Its first row query returns either
$\operatorname{HMAC}_h(K,m)$ or an independent uniform 256-bit tag; at most
$T\ge1$ row calls repeat this same transcript and cached answer. The algorithm
has no additional raw-key, midstate or key-dependent output interface.
It may continue querying the same $h$ and output a bit. The budget $q$ counts
all public compression queries, including preparation, other real rows,
continuation and the terminal test. Independently sampled known keys may
coincide with $K$; their ordinary calls to $h$ remain allowed and counted.

Let $\epsilon_{\rm row}(T,q,L)$ be any upper bound on this experiment's
bit-distinguishing gap. This operational definition is the primitive input to
the composition theorem. Appendix~\ref{app:primitive-inputs} gives the
compression-random-oracle instantiation, its padding convention and query
accounting. A zero answer is retained for the protocol's actual error check.

Define $H_k$, for $0\le k\le N$, by running the stopped experiment with every
unprotected row real, every protected row $i\le k$ real, and every protected
row $i>k$ replaced by its own independent uniform cached tag.  All later
recipes, flags, errors, and selections use that world's actual replies.
Thus $H_N=G^{\rm stop}$ and $H_0=:S^{\rm stop}$.  A zero ideal tag takes the
original error branch; it is never resampled.  An absent invocation makes no
row query.

\begin{theorem}[Reverse-order report composition]
\label{thm:report-reverse-composition}
For each slot $i$, let $L_i,T_i,q_i$ uniformly bound the transcript length,
row calls, and complete ordinary compression work of the reduction below,
including its prefix, future computation, other real rows, and final test.
Set
\[
 \epsilon_{{\rm rep},i}
   =\epsilon_{\rm row}(T_i,q_i,L_i),
 \qquad
 \epsilon_{{\rm rep},\Sigma}
   =\min\!\left\{1,\sum_{i=1}^N\epsilon_{{\rm rep},i}\right\}.
\]
Under the stated interface, every common efficient terminal bit test $D$
satisfies
\[
 \left|\Pr[G^{\rm stop}:D=1]
       -\Pr[S^{\rm stop}:D=1]\right|
 \le\epsilon_{{\rm rep},\Sigma}.
\]
For a terminal score $Z\in[-1,1]$ which is zero on stopping,
\[
 \left|\mathbb E_{G^{\rm stop}}Z
       -\mathbb E_{S^{\rm stop}}Z\right|
 \le2\epsilon_{{\rm rep},\Sigma}.
\]
The tests use the common available view, semantic ledger, protocol-interface
outputs, and independent test coins, including the retained application
challenge bit.  They have no extra access to protected key or transcript bytes.
\end{theorem}

\begin{proof}
Fix $i$ and compare $H_i$ with $H_{i-1}$.  Run the full prefix from setup,
locally sampling all earlier keys and computing their real rows.  Before the
selected row is used, both hybrids have the same prefix rules.  If slot $i$
is absent or unprotected, their behavior agrees and no row challenge is used.
Otherwise install the one-row challenger's key as an opaque
$\mathsf{KeyRef}(i)$.  The transcript $m_i$ is computable: it excludes its own
key, all referenced earlier keys were locally sampled, and the source atoms
are still available at this proof stage.  Submit this actual transcript to
the row oracle when its first MAC call is reached, retain its reply, and run
the original nonzero and persistence checks.

Maintain the original-report, key-origin, recipe, durable-state, and event
tables.  Every available output agrees with the selected hybrid.  Earlier
real rows remain computable because their immutable transcripts predate
$K_i$ and their keys are known.  A future unprotected row cannot depend on
$\mathsf{KeyRef}(i)$: such a dependency would make it protected by
\eqref{eq:report-protection}.  Compute that row with its locally sampled key
and public $h$.  A future protected row has index greater than $i$ and is
already ideal in both hybrids; sample its first tag, cache it by original
invocation, and retain its symbolic recipe without evaluating the unavailable
bytes.  This includes a later report which contains $K_i$ or another earlier
protected key and tag.  Flags and subsequent known fields are recomputed
using the actual returned tags.

For replay and validation, the selected row is called only on its original
$m_i$; preserve each counted recomputation and compare supplied identifiers
with its retained reply.  Copying, failure after a durable commit, and
crash/resume use the same immutable handles and the actual reached storage
state. Each report retains its original key and immutable transcript. Public decisions
use available metadata and common-interface outputs.  Before an outward
protected-key projection, the same ledger event
$\mathsf{Bad}_{\rm rep}$ enters the common stopped state with application
score zero; any allowed terminal-bit postprocessing uses that state.  No
unavailable key bytes are serialized.  Unprotected-key fragments retain their known bytes.
This gives an explicit online simulation throughout the continuation.

The selected real row therefore produces $H_i$, and its ideal counterpart
produces $H_{i-1}$.  All other real keys are sampled as actual independent byte
strings.  If one happens to equal the unknown challenged key, evaluating its
real MAC with public $h$ automatically preserves that equality; those queries
belong in $q_i$.  No comparison with the hidden key or distinct-key conditioning
is used.  Intentional aliases of the same original report instead share its
single row cache.

The reduction executes the adversary's quantum queries to the same $h$ and
introduces no measurement of those queries or oracle-dependent initial advice.
The one-row lemma bounds the adjacent bit difference by
$\epsilon_{{\rm rep},i}$.  Sum over the reverse-order hybrids and cap at one.
For the score, use the allowed randomized terminal test which outputs one
with probability $(1+Z)/2$; its difference is half the score-expectation
difference, giving the displayed factor two.
\end{proof}

\begin{corollary}[Original eligibility and real disclosures]
\label{cor:report-exposure-residual}
For the full real experiment $G$, define
\[
 \eta_{\rm rep}(Z)
 =\mathbb E_G\!\left[|Z|\,
                    \mathbf1_{\mathsf{Bad}_{\rm rep}}\right].
\]
Then
\[
 \left|\mathbb E_GZ-\mathbb E_{S^{\rm stop}}Z\right|
 \le\min\!\left\{2,\eta_{\rm rep}(Z)
                       +2\epsilon_{{\rm rep},\Sigma}\right\},
\]
and hence
\[
 |\mathbb E_GZ|
 \le\min\!\left\{1,\eta_{\rm rep}(Z)
       +2\epsilon_{{\rm rep},\Sigma}
       +|\mathbb E_{S^{\rm stop}}Z|\right\}.
\]
For $Z=\mathbf1_{\mathsf{Eligible}}(-1)^{b\oplus b'}$,
$\eta_{\rm rep}(Z)=
\Pr_G[\mathsf{Eligible}\land\mathsf{Bad}_{\rm rep}]$.
\end{corollary}

\begin{proof}
Couple $G$ and $G^{\rm stop}$ up to the first
$\mathsf{Bad}_{\rm rep}$.  On its complement the runs agree; on that event the
stopped score is zero while $G$ continues with the actual disclosed bytes.
Therefore
\[
 \mathbb E_GZ-\mathbb E_{G^{\rm stop}}Z
   =\mathbb E_G[Z\mathbf1_{\mathsf{Bad}_{\rm rep}}].
\]
Bound its magnitude by $\eta_{\rm rep}(Z)$ and apply
Theorem~\ref{thm:report-reverse-composition}.  No conditioning on eligibility
or independence of later target selection is used.
\end{proof}

\paragraph{Consequences for the record connection.}
The residual is zero on a disclosure surface whose actual provenance already
makes the original target permanently ineligible.  A batch-only report-key
fragment requires its own mapping: it is not automatically a member-session
capture.  This distinction is material.  Given a known transcript and its
later disclosed key, recomputing HMAC distinguishes the real tag from an
independent tag with difference $1-2^{-256}$.  This example does not assert
that a native batch fragment also reveals its members' complete transcript.

The report replacement precedes KEM embedding.  Afterwards every protected
row needs only its original cached tag, and every unprotected recipe is free
of $\mathcal X$ and protected raw-key dependencies.  The report layer can
therefore use opaque source handles for tag generation and immutable tag
validation.  Direct source-atom or application-plaintext outputs retain their
own adapters and acquisition rules.  The preceding private-record hop keeps
its own error and acquisition correspondence; this theorem does not infer
that correspondence from the report stopping game.  One precise native
composition first charges the actual eligible-exposure residual, then
abstracts the stopped native experiment, and finally applies the report
hybrids.  The record correspondence must be proved for that explicit stopped
interface; a key-only fragment does not inherit the original renderer's
permanent-target-failure premise.  Likewise, arbitrary
changed-byte import is a separate MAC-query interface from the immutable
validation proved here.  These are specific interface boundaries of the
composition, while the dependent immutable rows themselves are covered by
the explicit reduction above.
